% ---------------------------------------------------------------------
%  3. Experimental setup
% ---------------------------------------------------------------------
\section{Experimental setup}
\label{sec:setup}

\subsection{Problems and split}
The kit ships sixteen sample problems in the holdout's format. I froze a
development set \Sdev{} ($n=9$) for tuning and an evaluation set \Seval{} ($n=7$),
and report Part~Two science on \Sdev. \Sdev{} spans easy algebra, number-theoretic
identities, and hard olympiad/Putnam statements, so a single number never hides
the difficulty spread.

\subsection{The metric: substantive, not merely accepted}
I track two signals and never conflate them:
\begin{itemize}
  \item \textbf{Comparator pass} --- the Lean file is accepted by the open-source
        comparator. This is the mechanical quantity the grader scores.
  \item \textbf{Substantive} --- the pass also proves the \emph{actual} statement,
        surviving the integrity and anti-tautology guards
        (\S\ref{sec:mechanism}). A pass that only restates the goal does not
        count.
\end{itemize}
\textbf{Every headline number in this report is substantive.} This is a stricter
bar than the task requires, and I adopt it on purpose: it is the only count that
survives an independent re-run with a statement-locked comparator, and it is what
separates a real result from an inflated one.

\subsection{Caps}
The judge enforces \$1 and 8\,h per problem. A few development-matrix runs use
tighter measurement caps so a full sweep fits a small budget; the judge-caps
stress run in \S\ref{sec:results} uses the full \$1/8\,h. Because a repair turn
costs a fraction of a cent, the binding constraint on hard problems was
wall-clock, not the dollar cap --- which is why the ladder can afford a free
sweep, sampled candidates, and a combine step before the expensive residual stage.

\subsection{Part~Two conditions}
To ask whether collaboration beats solo, I hold the code path and caps fixed and
vary only the roles: \SQ/\SG{} (each model solo, on the kit's existing multi-turn
Lean-repair baseline), \RQ/\RG{} (explicit propose/repair, same model), and
\HQG/\HGQ{} (cross-model handoff, each direction). Two points of hygiene. The solo
conditions use the kit's real repair baseline --- I do \emph{not} redefine
``solo'' as a weaker single shot to flatter collaboration. And the \emph{Solo
Union} $\Uni = \SQ \lor \SG$ is \emph{derived}, not a run: it is the natural
non-collaborative baseline any coordination layer must beat to justify itself.

\subsection{Confounds, stated once}
So the results read as prose, the caveats live here: solo baselines sometimes had
more turns than the repair budgets (secondary metrics reported, not perfectly
matched); some matrix cells use measurement caps, not judge caps; with $n=9$ a
one-problem difference is within noise, so I read patterns, not decimals; cells
are single runs of stochastic models; and the full16 stress run is my own sample
suite, not the grader's private holdout.
